\documentclass[11pt, a4paper]{article}
\usepackage[a4paper, top=2.5cm, bottom=2.5cm, left=2cm, right=2cm]{geometry}
\usepackage{fontspec}
\usepackage[english, bidi=basic, provide=*]{babel}
\babelprovide[import, onchar=ids fonts]{english}
\babelfont{rm}{TeX Gyre Termes}
\babelfont{sf}{TeX Gyre Heros}
\babelfont{tt}{TeX Gyre Cursor}
\usepackage{enumitem}
\usepackage{amsmath}

\title{MCQ Practice Set: Meiosis and Sexual Life Cycles}
\author{Subject: Biology (Chapter 13)}
\date{}

\begin{document}
\maketitle

\section*{Practice Questions}

\begin{enumerate}[leftmargin=*, label=\arabic*.]

\item \textbf{The transmission of traits from one generation to the next is best defined as:}
\begin{enumerate}[label=\Alph*.]
\item Variation
\item Genetics
\item Heredity
\item Locus
\item Clone
\end{enumerate}
\textbf{Answer:} C \\
\textbf{Explanation:} Heredity, or inheritance, is the generation-to-generation transmission of traits.

\item \textbf{Which of the following is the specific location of a gene along the length of a chromosome?}
\begin{enumerate}[label=\Alph*.]
\item Allele
\item Locus
\item Centromere
\item Chiasma
\item Kinetochore
\end{enumerate}
\textbf{Answer:} B \\
\textbf{Explanation:} A locus is the precise physical location of a gene on a chromosome.

\item \textbf{In asexual reproduction, a single individual passes copies of all its genes to its offspring without:}
\begin{enumerate}[label=\Alph*.]
\item Mitosis
\item DNA replication
\item The fusion of gametes
\item Mutations
\item Growth
\end{enumerate}
\textbf{Answer:} C \\
\textbf{Explanation:} Asexual reproduction involves a single parent and does not involve the fusion of sperm and egg.

\item \textbf{A group of genetically identical individuals is called a:}
\begin{enumerate}[label=\Alph*.]
\item Species
\item Population
\item Clone
\item Homolog
\item Recombinant
\end{enumerate}
\textbf{Answer:} C \\
\textbf{Explanation:} Clones are individuals that are genetically identical to the parent, usually produced via mitosis.

\item \textbf{Which type of cell in the human body contains 46 chromosomes?}
\begin{enumerate}[label=\Alph*.]
\item Sperm
\item Egg
\item Gamete
\item Somatic cell
\item Both A and B
\end{enumerate}
\textbf{Answer:} D \\
\textbf{Explanation:} Somatic cells are all body cells except gametes and contain two sets of chromosomes.

\item \textbf{The two chromosomes of a pair that have the same length, centromere position, and staining pattern are:}
\begin{enumerate}[label=\Alph*.]
\item Sister chromatids
\item Recombinant chromosomes
\item Homologous chromosomes
\item Sex chromosomes
\item Autosomes
\end{enumerate}
\textbf{Answer:} C \\
\textbf{Explanation:} Homologous chromosomes (homologs) carry genes controlling the same inherited characters at equivalent loci.

\item \textbf{Human females have a homologous pair of \underline{\hspace{1cm}} chromosomes, while males have \underline{\hspace{1cm}}.}
\begin{enumerate}[label=\Alph*.]
\item XY; XX
\item XX; XY
\item YY; XX
\item XX; YY
\item X; Y
\end{enumerate}
\textbf{Answer:} B \\
\textbf{Explanation:} Females typically have XX; males have XY, where only small parts are homologous.

\item \textbf{Chromosomes that do not determine the sex of an individual are called:}
\begin{enumerate}[label=\Alph*.]
\item Sex chromosomes
\item Loci
\item Autosomes
\item Allosomes
\item Chromatids
\end{enumerate}
\textbf{Answer:} C \\
\textbf{Explanation:} Of the 23 pairs of human chromosomes, 22 pairs are autosomes.

\item \textbf{Any cell with two sets of chromosomes is called a \underline{\hspace{1cm}} cell.}
\begin{enumerate}[label=\Alph*.]
\item Haploid
\item Diploid
\item Zygote
\item Somatic
\item Polyploid
\end{enumerate}
\textbf{Answer:} B \\
\textbf{Explanation:} Diploid cells (2n) contain two chromosome sets, one from each parent.

\item \textbf{For humans, the haploid number ($n$) is \underline{\hspace{1cm}}.}
\begin{enumerate}[label=\Alph*.]
\item 46
\item 92
\item 23
\item 12
\item 8
\end{enumerate}
\textbf{Answer:} C \\
\textbf{Explanation:} The haploid set consists of 22 autosomes plus a single sex chromosome.

\item \textbf{A zygote is formed through the union of gametes in a process called:}
\begin{enumerate}[label=\Alph*.]
\item Meiosis
\item Mitosis
\item Fertilization
\item Budding
\item Synapsis
\end{enumerate}
\textbf{Answer:} C \\
\textbf{Explanation:} Fertilization is the fusion of haploid gametes to restore the diploid state.

\item \textbf{What process reduces the chromosome number from diploid to haploid?}
\begin{enumerate}[label=\Alph*.]
\item Mitosis
\item Binary fission
\item Meiosis
\item Fertilization
\item Replication
\end{enumerate}
\textbf{Answer:} C \\
\textbf{Explanation:} Meiosis is a special type of cell division that halves the chromosome set count.

\item \textbf{In plants and some algae, the life cycle includes both diploid and haploid multicellular stages called:}
\begin{enumerate}[label=\Alph*.]
\item Binary fission
\item Alternation of generations
\item Sexual dimorphism
\item Parthenogenesis
\item Cloning
\end{enumerate}
\textbf{Answer:} B \\
\textbf{Explanation:} This cycle involves a sporophyte (2n) and a gametophyte (n).

\item \textbf{The multicellular diploid stage in the plant life cycle is the:}
\begin{enumerate}[label=\Alph*.]
\item Gametophyte
\item Spore
\item Zygote
\item Sporophyte
\item Pollen
\end{enumerate}
\textbf{Answer:} D \\
\textbf{Explanation:} The sporophyte produces haploid spores via meiosis.

\item \textbf{Unlike animals, most fungi and some protists:}
\begin{enumerate}[label=\Alph*.]
\item Have no haploid stage
\item Form a multicellular diploid organism
\item Have a diploid zygote as the only diploid stage
\item Reproduce only asexually
\item Use mitosis to produce spores
\end{enumerate}
\textbf{Answer:} C \\
\textbf{Explanation:} In these organisms, the zygote undergoes meiosis immediately, producing haploid cells that divide by mitosis.

\item \textbf{Meiosis is preceded by \underline{\hspace{1cm}}, which includes \underline{\hspace{1cm}} phase.}
\begin{enumerate}[label=\Alph*.]
\item Prophase; S
\item Interphase; S
\item Cytokinesis; G1
\item Anaphase; G2
\item Telophase; M
\end{enumerate}
\textbf{Answer:} B \\
\textbf{Explanation:} DNA replication occurs during the S phase of interphase before meiosis starts.

\item \textbf{The first division of meiosis, Meiosis I, results in:}
\begin{enumerate}[label=\Alph*.]
\item Separation of sister chromatids
\item Four haploid cells
\item Separation of homologous chromosomes
\item Two diploid cells
\item DNA replication
\end{enumerate}
\textbf{Answer:} C \\
\textbf{Explanation:} Meiosis I is the reductional division where homologs move to opposite poles.

\item \textbf{What happens during synapsis in Prophase I?}
\begin{enumerate}[label=\Alph*.]
\item Sister chromatids separate
\item The nuclear envelope reforms
\item Homologs pair up and are held by a synaptonemal complex
\item DNA is replicated
\item Spindle fibers disappear
\end{enumerate}
\textbf{Answer:} C \\
\textbf{Explanation:} Synapsis is the state where homologs are physically connected along their lengths.

\item \textbf{The X-shaped regions where crossing over has occurred are called:}
\begin{enumerate}[label=\Alph*.]
\item Kinetochores
\item Centromeres
\item Chiasmata
\item Cohesins
\item Asters
\end{enumerate}
\textbf{Answer:} C \\
\textbf{Explanation:} Chiasmata are the physical manifestations of crossing over between nonsister chromatids.

\item \textbf{At Metaphase I, what is aligned at the metaphase plate?}
\begin{enumerate}[label=\Alph*.]
\item Individual chromosomes
\item Pairs of homologous chromosomes
\item Sister chromatids
\item Single-stranded DNA
\item Nucleoli
\end{enumerate}
\textbf{Answer:} B \\
\textbf{Explanation:} Pairs of homologs line up, with one chromosome of each pair facing each pole.

\item \textbf{During Anaphase I, what allows the homologs to separate?}
\begin{enumerate}[label=\Alph*.]
\item Cleavage of cohesins along the chromatid arms
\item Cleavage of cohesins at the centromeres
\item Depolymerization of all microtubules
\item Reappearance of the nucleolus
\item Fusion of sister chromatids
\end{enumerate}
\textbf{Answer:} A \\
\textbf{Explanation:} Cohesion along arms is released, but centromere cohesion remains to keep sister chromatids together.

\item \textbf{Meiosis II is very similar to mitosis because:}
\begin{enumerate}[label=\Alph*.]
\item DNA is replicated again
\item Homologs pair up
\item Sister chromatids separate
\item It results in diploid cells
\item Synapsis occurs
\end{enumerate}
\textbf{Answer:} C \\
\textbf{Explanation:} In Meiosis II (the equational division), sister chromatids are pulled apart.

\item \textbf{Which protein complex holds sister chromatids together?}
\begin{enumerate}[label=\Alph*.]
\item Tubulin
\item Actin
\item Cohesin
\item Myosin
\item Separase
\end{enumerate}
\textbf{Answer:} C \\
\textbf{Explanation:} Cohesins mediate the attachment of sister chromatids until they are cleaved.

\item \textbf{What is the total number of combinations possible when chromosomes sort independently in humans?}
\begin{enumerate}[label=\Alph*.]
\item 46
\item 23
\item $2^{23}$ ($\approx 8.4$ million)
\item $2^{46}$
\item 70 trillion
\end{enumerate}
\textbf{Answer:} C \\
\textbf{Explanation:} The number of combinations is $2^n$, where $n=23$ for humans.

\item \textbf{Crossing over produces \underline{\hspace{1cm}} chromosomes, which carry genes from two different parents.}
\begin{enumerate}[label=\Alph*.]
\item Autosomal
\item Homologous
\item Recombinant
\item Sister
\item Parental
\end{enumerate}
\textbf{Answer:} C \\
\textbf{Explanation:} Recombinant chromosomes are individual chromosomes that carry a mix of maternal and paternal DNA.

\item \textbf{In meiosis, how many times is the DNA replicated and how many times does the cell divide?}
\begin{enumerate}[label=\Alph*.]
\item Twice; twice
\item Once; once
\item Once; twice
\item Twice; once
\item Zero; twice
\end{enumerate}
\textbf{Answer:} C \\
\textbf{Explanation:} DNA is replicated once in interphase, followed by two rounds of division (I and II).

\item \textbf{The zipper-like structure that holds homologs together is the:}
\begin{enumerate}[label=\Alph*.]
\item Centrosome
\item Synaptonemal complex
\item Kinetochore
\item Spindle apparatus
\item Aster
\end{enumerate}
\textbf{Answer:} B \\
\textbf{Explanation:} The synaptonemal complex facilitates synapsis and crossing over.

\item \textbf{Which mechanism does NOT contribute to genetic variation in sexual reproduction?}
\begin{enumerate}[label=\Alph*.]
\item Independent assortment
\item Crossing over
\item Random fertilization
\item Mitosis
\item Mutation
\end{enumerate}
\textbf{Answer:} D \\
\textbf{Explanation:} Mitosis conserves genetic identity; it does not generate variation.

\item \textbf{Natural selection results in the accumulation of genetic variations favored by the \underline{\hspace{1cm}}.}
\begin{enumerate}[label=\Alph*.]
\item Nucleus
\item Environment
\item Mitochondria
\item Golgi apparatus
\item Individual
\end{enumerate}
\textbf{Answer:} B \\
\textbf{Explanation:} Individuals best suited to the local environment leave more offspring.

\item \textbf{Which organism is an exception to the rule that sex is needed for genetic diversity?}
\begin{enumerate}[label=\Alph*.]
\item Human
\item Pea plant
\item Bdelloid rotifer
\item Yeast
\item Hydra
\end{enumerate}
\textbf{Answer:} C \\
\textbf{Explanation:} Bdelloids reproduce asexually but pick up foreign DNA to maintain diversity.

\item \textbf{A human cell containing 22 autosomes and a Y chromosome is a:}
\begin{enumerate}[label=\Alph*.]
\item Somatic cell
\item Zygote
\item Sperm
\item Egg
\item Diploid cell
\end{enumerate}
\textbf{Answer:} C \\
\textbf{Explanation:} This is a haploid gamete containing the male sex chromosome.

\item \textbf{Homologous chromosomes move toward opposite poles during:}
\begin{enumerate}[label=\Alph*.]
\item Mitosis
\item Meiosis I
\item Meiosis II
\item Fertilization
\item Interphase
\end{enumerate}
\textbf{Answer:} B \\
\textbf{Explanation:} This occurs specifically during Anaphase I.

\item \textbf{The DNA content of a diploid cell in $G_1$ is $x$. What is the DNA content at Metaphase I?}
\begin{enumerate}[label=\Alph*.]
\item $0.5x$
\item $x$
\item $2x$
\item $4x$
\item $0.25x$
\end{enumerate}
\textbf{Answer:} C \\
\textbf{Explanation:} DNA is replicated in S phase, doubling the content from $x$ to $2x$.

\item \textbf{Following the cell in the previous question, what is the DNA content in a single cell at Metaphase II?}
\begin{enumerate}[label=\Alph*.]
\item $0.5x$
\item $x$
\item $2x$
\item $4x$
\item $0.25x$
\end{enumerate}
\textbf{Answer:} B \\
\textbf{Explanation:} After Meiosis I, the content is halved (from $2x$ to $x$).

\item \textbf{Which stage is characterized by the sudden parting of sister chromatids in meiosis?}
\begin{enumerate}[label=\Alph*.]
\item Anaphase I
\item Anaphase II
\item Metaphase I
\item Prophase II
\item Telophase I
\end{enumerate}
\textbf{Answer:} B \\
\textbf{Explanation:} Sister chromatids separate during Anaphase II.

\item \textbf{In humans, meiosis occurs only in:}
\begin{enumerate}[label=\Alph*.]
\item Skin cells
\item Liver cells
\item Germ cells in the gonads
\item Bone marrow
\item Nerve cells
\end{enumerate}
\textbf{Answer:} C \\
\textbf{Explanation:} Only specialized cells in the ovaries or testes undergo meiosis.

\item \textbf{A version of a gene at a specific locus is called a/an:}
\begin{enumerate}[label=\Alph*.]
\item Mutation
\item Allele
\item Homolog
\item Chromosome
\item Clone
\end{enumerate}
\textbf{Answer:} B \\
\textbf{Explanation:} Alleles are different versions of a gene that may exist at a locus.

\item \textbf{Crossing over begins during which phase?}
\begin{enumerate}[label=\Alph*.]
\item Prophase I
\item Metaphase I
\item Prophase II
\item Anaphase I
\item Interphase
\end{enumerate}
\textbf{Answer:} A \\
\textbf{Explanation:} Crossing over occurs during the early stages of Prophase I.

\item \textbf{The number of sets of chromosomes in a gametophyte is \underline{\hspace{1cm}}.}
\begin{enumerate}[label=\Alph*.]
\item 1
\item 2
\item 3
\item 4
\item 0
\end{enumerate}
\textbf{Answer:} A \\
\textbf{Explanation:} The gametophyte is the haploid multicellular stage.

\item \textbf{Which of the following describes a "reductional" division?}
\begin{enumerate}[label=\Alph*.]
\item Mitosis
\item Meiosis I
\item Meiosis II
\item DNA synthesis
\item Cytokinesis
\end{enumerate}
\textbf{Answer:} B \\
\textbf{Explanation:} Meiosis I reduces the number of chromosome sets from two to one.

\item \textbf{How many chromatids are in a human cell at Metaphase I?}
\begin{enumerate}[label=\Alph*.]
\item 23
\item 46
\item 92
\item 184
\item 12
\end{enumerate}
\textbf{Answer:} C \\
\textbf{Explanation:} 46 duplicated chromosomes each have 2 chromatids ($46 \times 2 = 92$).

\item \textbf{Sister chromatid cohesion is released at centromeres during:}
\begin{enumerate}[label=\Alph*.]
\item Anaphase I
\item Anaphase II
\item Metaphase I
\item Metaphase II
\item Telophase II
\end{enumerate}
\textbf{Answer:} B \\
\textbf{Explanation:} Release at the centromere allows chromatids to move to opposite poles.

\item \textbf{What is the evolutionary source of different alleles?}
\begin{enumerate}[label=\Alph*.]
\item Natural selection
\item Crossing over
\item Mutation
\item Independent assortment
\item Fertilization
\end{enumerate}
\textbf{Answer:} C \\
\textbf{Explanation:} Mutations (changes in DNA) are the original source of genetic diversity.

\item \textbf{Which phase follows Meiosis I with no further DNA replication?}
\begin{enumerate}[label=\Alph*.]
\item Prophase I
\item S phase
\item Interkinesis/Prophase II
\item G2 phase
\item Mitosis
\end{enumerate}
\textbf{Answer:} C \\
\textbf{Explanation:} There is no S phase between Meiosis I and Meiosis II.

\item \textbf{If $n=3$, how many possible combinations can result from independent assortment?}
\begin{enumerate}[label=\Alph*.]
\item 3
\item 6
\item 8
\item 9
\item 12
\end{enumerate}
\textbf{Answer:} C \\
\textbf{Explanation:} $2^3 = 8$.

\item \textbf{Maternal and paternal chromosomes in a pair are called \underline{\hspace{1cm}}.}
\begin{enumerate}[label=\Alph*.]
\item Sister chromatids
\item Nonsister chromatids
\item Homologs
\item Both B and C
\item Clones
\end{enumerate}
\textbf{Answer:} D \\
\textbf{Explanation:} They are homologs, and the chromatids between them are nonsister.

\item \textbf{A display of condensed chromosomes arranged in pairs is a:}
\begin{enumerate}[label=\Alph*.]
\item Genotype
\item Phenotype
\item Karyotype
\item Locus
\item Pedigree
\end{enumerate}
\textbf{Answer:} C \\
\textbf{Explanation:} Karyotypes show chromosome pairs ordered by size.

\item \textbf{In which phase do the nuclear envelopes reform and cytokinesis occurs to produce four cells?}
\begin{enumerate}[label=\Alph*.]
\item Telophase I
\item Telophase II
\item Anaphase II
\item Metaphase II
\item Prophase II
\end{enumerate}
\textbf{Answer:} B \\
\textbf{Explanation:} Telophase II and cytokinesis finalize the production of four haploid cells.

\item \textbf{Cohesion along the chromatid arms is released at the start of:}
\begin{enumerate}[label=\Alph*.]
\item Anaphase I
\item Anaphase II
\item Metaphase I
\item Metaphase II
\item Prophase I
\end{enumerate}
\textbf{Answer:} A \\
\textbf{Explanation:} This release allows the homologous chromosomes to part ways.

\item \textbf{What is the primary significance of the bdelloid rotifer state of "suspended animation"?}
\begin{enumerate}[label=\Alph*.]
\item It prevents mutations
\item It allows for sexual reproduction
\item It lets foreign DNA enter via cracked membranes
\item It increases the diploid number
\item It stops all evolution
\end{enumerate}
\textbf{Answer:} C \\
\textbf{Explanation:} This provides a non-sexual mechanism for achieving genetic diversity.

\item \textbf{A "recombinant" chromosome contains DNA from:}
\begin{enumerate}[label=\Alph*.]
\item Only one parent
\item Two different parents
\item The mitochondria only
\item The ribosomes only
\item Mutations only
\end{enumerate}
\textbf{Answer:} B \\
\textbf{Explanation:} It results from the exchange of segments between maternal and paternal homologs.

\item \textbf{Which is the "equational" division?}
\begin{enumerate}[label=\Alph*.]
\item Meiosis I
\item Meiosis II
\item Fertilization
\item Binary fission
\item DNA replication
\end{enumerate}
\textbf{Answer:} B \\
\textbf{Explanation:} Meiosis II is equational because it separates sister chromatids, maintaining the haploid state.

\item \textbf{Spores in plants are produced by which stage?}
\begin{enumerate}[label=\Alph*.]
\item Gametophyte
\item Sporophyte
\item Pollen tube
\item Endosperm
\item Zygote
\end{enumerate}
\textbf{Answer:} B \\
\textbf{Explanation:} The sporophyte undergoes meiosis to produce haploid spores.

\item \textbf{Which event occurs in Prophase I but NOT in mitotic prophase?}
\begin{enumerate}[label=\Alph*.]
\item Condensation of chromosomes
\item Spindle formation
\item Synapsis and crossing over
\item Disappearance of the nucleolus
\item Fragmentation of the nuclear envelope
\end{enumerate}
\textbf{Answer:} C \\
\textbf{Explanation:} Synapsis is unique to meiosis.

\item \textbf{How many daughter cells are produced at the end of meiosis?}
\begin{enumerate}[label=\Alph*.]
\item 1
\item 2
\item 3
\item 4
\item 8
\end{enumerate}
\textbf{Answer:} D \\
\textbf{Explanation:} Meiosis produces four genetically distinct haploid daughter cells.

\item \textbf{Independent assortment occurs because of the random orientation of:}
\begin{enumerate}[label=\Alph*.]
\item Sister chromatids at Metaphase II
\item Homologous pairs at Metaphase I
\item Spindle poles in Prophase I
\item Genes on a chromosome
\item DNA during S phase
\end{enumerate}
\textbf{Answer:} B \\
\textbf{Explanation:} Each pair of homologs positions itself independently of other pairs.

\item \textbf{What is the probability that a human sperm gets the maternal chromosome of a specific pair?}
\begin{enumerate}[label=\Alph*.]
\item 100\%
\item 0\%
\item 50\%
\item 25\%
\item 75\%
\end{enumerate}
\textbf{Answer:} C \\
\textbf{Explanation:} Like a coin flip, there is a 50\% chance for either homolog.

\item \textbf{The fusion of human gametes results in \underline{\hspace{1cm}} possible diploid combinations.}
\begin{enumerate}[label=\Alph*.]
\item 8.4 million
\item 46
\item 70 trillion
\item 23
\item Infinite
\end{enumerate}
\textbf{Answer:} C \\
\textbf{Explanation:} $2^{23} \times 2^{23}$ results in about 70 trillion combinations.

\item \textbf{At the end of Anaphase I, the cells are \underline{\hspace{1cm}}.}
\begin{enumerate}[label=\Alph*.]
\item Diploid
\item Haploid
\item Tetrads
\item Somatic
\item Clones
\end{enumerate}
\textbf{Answer:} B \\
\textbf{Explanation:} The chromosome number is halved by the end of Meiosis I.

\item \textbf{Chromosomes are duplicated \underline{\hspace{1cm}}.}
\begin{enumerate}[label=\Alph*.]
\item Twice during meiosis
\item Once before Meiosis I only
\item Once before Meiosis II only
\item Once before both Meiosis I and II
\item Not at all
\end{enumerate}
\textbf{Answer:} B \\
\textbf{Explanation:} Duplication only occurs once, during the S phase of interphase.

\item \textbf{An unfertilized human egg contains \underline{\hspace{1cm}} sex chromosome(s).}
\begin{enumerate}[label=\Alph*.]
\item One X
\item One Y
\item One X or one Y
\item Two X
\item Zero
\end{enumerate}
\textbf{Answer:} A \\
\textbf{Explanation:} Eggs always carry an X chromosome.

\item \textbf{What causes offspring to resemble parents more than unrelated individuals?}
\begin{enumerate}[label=\Alph*.]
\item Shared environment
\item Inheritance of genes
\item Same diet
\item Identical mutations
\item Random chance
\end{enumerate}
\textbf{Answer:} B \\
\textbf{Explanation:} Genes are the hereditary units transmitted from parents to offspring.

\item \textbf{In meiosis II, the sister chromatids are \underline{\hspace{1cm}}.}
\begin{enumerate}[label=\Alph*.]
\item Identical
\item Not genetically identical (due to crossing over)
\item Replicated again
\item Not separated
\item Fused together
\end{enumerate}
\textbf{Answer:} B \\
\textbf{Explanation:} Crossing over makes the sister chromatids distinct.

\item \textbf{Bdelloid rotifers survive environmental stress by entering:}
\begin{enumerate}[label=\Alph*.]
\item Mitosis
\item Meiosis
\item Suspended animation
\item Fertilization
\item Binary fission
\end{enumerate}
\textbf{Answer:} C \\
\textbf{Explanation:} This state allows them to survive drying and take in foreign DNA.

\item \textbf{If an organism has $2n = 10$, its gametes have \underline{\hspace{1cm}} chromosomes.}
\begin{enumerate}[label=\Alph*.]
\item 2
\item 5
\item 10
\item 20
\item 1
\end{enumerate}
\textbf{Answer:} B \\
\textbf{Explanation:} $n = 2n / 2 = 5$.

\item \textbf{Which stage is the "first meiotic division"?}
\begin{enumerate}[label=\Alph*.]
\item Prophase I
\item Meiosis I
\item Meiosis II
\item Cytokinesis
\item S phase
\end{enumerate}
\textbf{Answer:} B \\
\textbf{Explanation:} Meiosis I is the first of two consecutive divisions.

\item \textbf{In Prophase I, DNA breaks are closed by joining to:}
\begin{enumerate}[label=\Alph*.]
\item The same chromatid
\item The sister chromatid
\item A nonsister chromatid
\item A different chromosome pair
\item The nuclear envelope
\end{enumerate}
\textbf{Answer:} C \\
\textbf{Explanation:} This joining creates the crossover between nonsister chromatids.

\item \textbf{The number of chiasmata in a human chromosome pair averages:}
\begin{enumerate}[label=\Alph*.]
\item 0
\item 1 to 3
\item 10 to 20
\item 46
\item 100
\end{enumerate}
\textbf{Answer:} B \\
\textbf{Explanation:} Typically, 1 to 3 crossover events occur per chromosome pair.

\item \textbf{Which of these occurs in both mitosis and Meiosis II?}
\begin{enumerate}[label=\Alph*.]
\item Synapsis
\item Reduction of chromosome sets
\item Separation of sister chromatids
\item Crossing over
\item Production of spores
\end{enumerate}
\textbf{Answer:} C \\
\textbf{Explanation:} Both processes are equational divisions regarding chromatids.

\item \textbf{What is the chromosome count in a human somatic cell after mitosis?}
\begin{enumerate}[label=\Alph*.]
\item 23
\item 46
\item 92
\item 12
\item 8
\end{enumerate}
\textbf{Answer:} B \\
\textbf{Explanation:} Mitosis conserves the diploid number of 46.

\item \textbf{In alternation of generations, spores divide by \underline{\hspace{1cm}} to form the gametophyte.}
\begin{enumerate}[label=\Alph*.]
\item Meiosis
\item Mitosis
\item Fertilization
\item Binary fission
\item Budding
\end{enumerate}
\textbf{Answer:} B \\
\textbf{Explanation:} Spores are already haploid and grow into a haploid stage via mitosis.

\item \textbf{What holds homologs together as the spindle forms in Meiosis I?}
\begin{enumerate}[label=\Alph*.]
\item Centromeres
\item Chiasmata and arm cohesion
\item Synaptonemal complex only
\item Kinetochores only
\item Microtubules
\end{enumerate}
\textbf{Answer:} B \\
\textbf{Explanation:} Chiasmata, resulting from crossing over, keep pairs together.

\item \textbf{The study of heredity and inherited variation is:}
\begin{enumerate}[label=\Alph*.]
\item Cytology
\item Genetics
\item Metabolism
\item Evolution
\item Anatomy
\end{enumerate}
\textbf{Answer:} B \\
\textbf{Explanation:} This is the formal definition of genetics.

\item \textbf{Haploid cells \underline{\hspace{1cm}} undergo meiosis.}
\begin{enumerate}[label=\Alph*.]
\item Always
\item Often
\item Cannot
\item Must
\item Only during stress
\end{enumerate}
\textbf{Answer:} C \\
\textbf{Explanation:} You cannot reduce a single set of chromosomes any further.

\item \textbf{The "waist" of a duplicated chromosome is the:}
\begin{enumerate}[label=\Alph*.]
\item Locus
\item Chiasma
\item Centromere
\item Chromatid
\item Aster
\end{enumerate}
\textbf{Answer:} C \\
\textbf{Explanation:} This is where the sister chromatids are most closely attached.

\item \textbf{Human sex cells have how many autosomes?}
\begin{enumerate}[label=\Alph*.]
\item 46
\item 23
\item 22
\item 1
\item 0
\end{enumerate}
\textbf{Answer:} C \\
\textbf{Explanation:} $n = 23$, which consists of 22 autosomes and 1 sex chromosome.

\item \textbf{Which is the source of unique combinations of genes?}
\begin{enumerate}[label=\Alph*.]
\item Asexual reproduction
\item Mutation only
\item Sexual reproduction
\item Binary fission
\item Mitotic cloning
\end{enumerate}
\textbf{Answer:} C \\
\textbf{Explanation:} Sexual reproduction reshuffles alleles to create uniqueness.

\item \textbf{How many chromatid arms does a duplicated chromosome have in Figure 13.4?}
\begin{enumerate}[label=\Alph*.]
\item 1
\item 2
\item 4
\item 8
\item 6
\end{enumerate}
\textbf{Answer:} C \\
\textbf{Explanation:} Each of the two sister chromatids has two arms ($2 \times 2 = 4$).

\item \textbf{The synaptonemal complex disassembles in \underline{\hspace{1cm}}.}
\begin{enumerate}[label=\Alph*.]
\item Prophase I
\item Metaphase I
\item Anaphase I
\item Telophase I
\item Meiosis II
\end{enumerate}
\textbf{Answer:} A \\
\textbf{Explanation:} It disassembles later in Prophase I, leaving chiasmata visible.

\item \textbf{Natural selection works by accumulating variations that are \underline{\hspace{1cm}}.}
\begin{enumerate}[label=\Alph*.]
\item Random
\item Dangerous
\item Advantageous
\item Identical
\item Asexual
\end{enumerate}
\textbf{Answer:} C \\
\textbf{Explanation:} Advantageous traits lead to higher reproductive success.

\end{enumerate}
\end{document}